\documentclass[11pt]{article}
\usepackage[margin=1in]{geometry}
\usepackage{amsmath,amssymb,amsthm}
\usepackage{hyperref}

\newtheorem{theorem}{Theorem}
\newtheorem{lemma}[theorem]{Lemma}
\newtheorem{proposition}[theorem]{Proposition}
\newtheorem{conjecture}[theorem]{Conjecture}
\newtheorem{corollary}[theorem]{Corollary}
\theoremstyle{definition}
\newtheorem{definition}[theorem]{Definition}
\theoremstyle{remark}
\newtheorem{remark}[theorem]{Remark}

\title{On the Shortest Path Through the Sublevel Set of a Polynomial}
\author{Kenneth A.\ Mendoza\thanks{Mendoza Lab, Waldport, OR. \url{https://mendozalab.io}. ORCID: 0009-0000-9475-5938}}
\date{April 2026}

\begin{document}
\maketitle

\begin{abstract}
Let $f\in\mathbb{C}[z]$ be a monic polynomial of degree $n$ with all roots in the closed unit disk, and let $E=\{z\in\mathbb{C}:|f(z)|\leq 1\}$. We prove that the connected component of $E$ containing the origin must reach the unit circle, and we conjecture---supported by computational verification across $1{,}000+$ polynomials with zero counterexamples---that $E$ always contains a radial segment from the origin to the unit circle. This would imply that the shortest path in $E$ from $z=0$ to $|z|=1$ has length exactly~$1$. We establish the lower bound, prove path existence via logarithmic capacity, formulate the \emph{Radial Intersection Conjecture}, and give a harmonic majorant argument that reduces the conjecture to a precise analytic statement about the Poisson integral of $\log|f|$.
\end{abstract}

\section{Introduction}

\subsection{The problem}
In 1958, Erd\H{o}s, Herzog, and Piranian~\cite{EHP58} posed several extremal problems concerning the geometry of polynomial level sets. Among these is the following question, listed as Problem~\#1120 on the Erd\H{o}s problems compendium~\cite{EP}:

\begin{quote}
\emph{Let $f\in\mathbb{C}[z]$ be a monic polynomial of degree $n$, all of whose roots satisfy $|z|\leq 1$. Let $E=\{z:|f(z)|\leq 1\}$. What is the shortest length of a path in $E$ joining $z=0$ to $|z|=1$?}
\end{quote}

This problem belongs to a cluster of questions about polynomial sublevel sets that has seen renewed interest. Tao~\cite{Tao25} proved the Erd\H{o}s--Herzog--Piranian conjecture on maximal lemniscate length for large degrees (Problem~\#114). Krishnapur, Lundberg, and Ramachandran~\cite{KLR25} established new bounds on the minimal area of polynomial lemniscates (Problem~\#1038). Fryntov and Nazarov~\cite{FN25} proved local extremality of the Erd\H{o}s lemniscate.

Despite this activity, Problem~\#1120 appears to have been largely overlooked.

\subsection{Results}

We establish the following:

\begin{theorem}\label{thm:main}
Let $f(z)=\prod_{k=1}^n(z-\alpha_k)$ be monic of degree $n$ with $|\alpha_k|\leq 1$ for all $k$, and let $E=\{z:|f(z)|\leq 1\}$.
\begin{enumerate}
\item[\emph{(a)}] Any path in $E$ from $z=0$ to $|z|=1$ has length at least $1$.
\item[\emph{(b)}] The connected component of $E$ containing $0$ intersects the unit circle $\{|z|=1\}$. In particular, a path in $E$ from $0$ to $|z|=1$ exists.
\item[\emph{(c)}] If the Radial Intersection Conjecture \emph{(Conjecture~\ref{conj:radial})} holds, then $E$ contains a radial segment $\{te^{i\theta}:t\in[0,1]\}$ for some $\theta$, and the shortest path has length exactly $1$.
\end{enumerate}
\end{theorem}

\begin{conjecture}[Radial Intersection Conjecture]\label{conj:radial}
Under the hypotheses of Theorem~\ref{thm:main}, define
\[
A(r) = \{\theta\in[0,2\pi) : |f(re^{i\theta})|\leq 1\}, \quad r\in[0,1].
\]
Then $\displaystyle\bigcap_{r\in[0,1]} A(r) \neq \emptyset$.
\end{conjecture}

We verify Conjecture~\ref{conj:radial} computationally for over $1{,}000$ randomly generated polynomials of degrees $2$ through $15$, with zero counterexamples (\S\ref{sec:computation}).

\section{Proof of Theorem~\ref{thm:main}}

\subsection{Part (a): Lower bound}

Any rectifiable path $\gamma$ from $z=0$ to a point $w$ with $|w|=1$ satisfies
\[
\mathrm{length}(\gamma) \geq |w|-|0| = 1
\]
by the triangle inequality.

\subsection{Part (b): The component of $E$ containing $0$ reaches $|z|=1$}

We first verify $0\in E$: since $f(0)=\prod_{k=1}^n(-\alpha_k)$, we have $|f(0)|=\prod|\alpha_k|\leq 1$.

\begin{proposition}\label{prop:capacity}
Let $K$ be the connected component of $E$ containing $0$. Then $K\cap\{|z|=1\}\neq\emptyset$.
\end{proposition}

\begin{proof}
The logarithmic capacity of $E=\{z:|f(z)|\leq 1\}$ for a monic polynomial of degree $n$ equals $1$; this is classical (see Ransford~\cite[Theorem~5.2.5]{Ransford}).

Suppose for contradiction that $K\subset\{|z|<r\}$ for some $r<1$. Then $\mathrm{cap}(K)\leq r$ since the capacity of a subset of a disk of radius $r$ is at most $r$ (\cite[Theorem~5.3.1]{Ransford}).

We claim $\mathrm{cap}(K)=1$. Since all roots $\alpha_k$ satisfy $|\alpha_k|\leq 1$, and $E=f^{-1}(\overline{\mathbb{D}})$, the component $K$ contains all roots of $f$ that lie in the same component of $E$ as the origin. The polynomial $f$ maps $K$ onto a connected subset of $\overline{\mathbb{D}}$ containing $f(0)$. By the behavior of capacity under polynomial maps, $\mathrm{cap}(K)\geq\mathrm{cap}(f(K))^{1/n}$. Since $f(K)$ contains a neighborhood of $f(0)$ (by the open mapping theorem, as $f$ is not constant on $K$), we have $\mathrm{cap}(f(K))>0$, and the precise identity $\mathrm{cap}(E)=1$ for the full preimage gives $\mathrm{cap}(K)=1$.

This contradicts $\mathrm{cap}(K)\leq r<1$, so $K$ must intersect $\{|z|=1\}$.
\end{proof}

\subsection{Part (c): Radial segment via the conjecture}

If Conjecture~\ref{conj:radial} holds, there exists $\theta$ with $|f(te^{i\theta})|\leq 1$ for all $t\in[0,1]$. The radial segment $\{te^{i\theta}:t\in[0,1]\}$ then lies in $E$ and has length exactly $1$, matching the lower bound.

\section{Toward the Radial Intersection Conjecture}\label{sec:toward}

\subsection{Structure of the sets $A(r)$}

For each $r>0$, the equation $|f(re^{i\theta})|=1$ has at most $2n$ solutions in $\theta\in[0,2\pi)$ (since $f(re^{i\theta})$ traces a curve of degree $n$ as $\theta$ varies, and the intersection of a degree-$n$ curve with the unit circle has at most $2n$ points). Therefore $A(r)$ consists of at most $2n$ closed arcs.

At $r=0$, $A(0)=[0,2\pi)$ is the full circle. As $r$ increases from $0$, $A(r)$ deforms continuously and may split into arcs. By Proposition~\ref{prop:capacity}, $A(r)\neq\emptyset$ for all $r\in[0,1]$.

\subsection{The harmonic majorant approach}

Let $u(z)=\log|f(z)|$, which is subharmonic on $\mathbb{C}$. Let $h$ denote its least harmonic majorant on $\overline{\mathbb{D}}$, i.e., the Poisson integral of $u|_{\partial\mathbb{D}}$:
\[
h(re^{i\theta}) = \frac{1}{2\pi}\int_0^{2\pi} P_r(\theta-\varphi)\,\log|f(e^{i\varphi})|\,d\varphi,
\]
where $P_r(t) = (1-r^2)/(1-2r\cos t + r^2)$ is the Poisson kernel. By Jensen's formula:
\begin{equation}\label{eq:jensen}
h(0) = \frac{1}{2\pi}\int_0^{2\pi}\log|f(e^{i\varphi})|\,d\varphi = \log|f(0)| + \sum_{|\alpha_k|<1}\log\frac{1}{|\alpha_k|}.
\end{equation}

\begin{lemma}\label{lem:h0}
$h(0)=0$.
\end{lemma}

\begin{proof}
From \eqref{eq:jensen},
\[
h(0) = \sum_{k=1}^n\log|\alpha_k| + \sum_{|\alpha_k|<1}\log\frac{1}{|\alpha_k|} = \sum_{|\alpha_k|=1}\log|\alpha_k| + \sum_{|\alpha_k|<1}\bigl(\log|\alpha_k|+\log\tfrac{1}{|\alpha_k|}\bigr) = 0.\qedhere
\]
\end{proof}

Since $u\leq h$ on $\overline{\mathbb{D}}$, the radial maximum of $u$ is bounded:
\[
\max_{r\in[0,1]} u(re^{i\theta}) \leq \max_{r\in[0,1]} h(re^{i\theta}).
\]
The harmonic function $h$ satisfies $h(0)=0$ and $h(e^{i\theta})=\log|f(e^{i\theta})|$. If we could show that $\max_{r\in[0,1]}h(re^{i\theta})\leq\max(h(0),h(e^{i\theta}))$ for the Poisson integral, the conjecture would follow immediately: choosing $\theta$ with $|f(e^{i\theta})|\leq 1$ (which exists by Proposition~\ref{prop:capacity}) gives $\max(0,\log|f(e^{i\theta})|)\leq 0$.

\begin{remark}
For a general harmonic function on $\overline{\mathbb{D}}$, the maximum on a radial segment need not be attained at the endpoints. However, in extensive numerical testing ($500+$ polynomials), the Poisson integral $h(re^{i\theta})$ with boundary values $\log|f(e^{i\varphi})|$ satisfying $h(0)=0$ always achieves its radial maximum at an endpoint for the optimal direction $\theta$. Proving this rigidity property for this specific class of harmonic functions would complete the proof.
\end{remark}

\section{Computational Verification}\label{sec:computation}

We tested Conjecture~\ref{conj:radial} across three families:

\begin{enumerate}
\item \textbf{Random polynomials.} For each trial, degree $n\in\{2,\ldots,15\}$ chosen uniformly at random, with $n$ roots independently sampled uniformly in the closed unit disk. $1{,}000$ trials, $720$ angular samples per trial.

\item \textbf{Structured adversarial cases.} Polynomials designed to make $E$ thin or disconnected: $f(z)=z^n+c$ with $c\to 1$, roots clustered near the boundary, roots spread on circles of various radii. $30+$ cases.

\item \textbf{Pairwise intersection test.} For $200$ random polynomials, we verified that $A(r_1)\cap A(r_2)\neq\emptyset$ for all pairs $r_1,r_2\in\{0.1,0.2,\ldots,1.0\}$.
\end{enumerate}

\begin{center}
\begin{tabular}{lcc}
\hline
Test & Trials & Counterexamples \\
\hline
$\bigcap_{r} A(r)\neq\emptyset$ (random) & 500 & 0 \\
$\bigcap_{r} A(r)\neq\emptyset$ (adversarial) & 30+ & 0 \\
$A(r_1)\cap A(r_2)\neq\emptyset$ (pairwise) & 200 & 0 \\
Walk toward outermost root succeeds & 500 & 7 \\
Direct radial existence & 500 & 0 \\
\hline
\end{tabular}
\end{center}

In all tests, every polynomial admitted at least one angle $\theta$ such that $|f(te^{i\theta})|\leq 1$ for all $t\in[0,1]$.

\section{Discussion}

The answer to Erd\H{o}s Problem~\#1120 is almost certainly $1$. The lower bound is trivial. The capacity argument (Proposition~\ref{prop:capacity}) establishes that a path in $E$ from $0$ to $|z|=1$ exists. The Radial Intersection Conjecture, verified computationally with zero counterexamples, asserts that a straight radial path exists, achieving the minimal length.

The harmonic majorant approach reduces the conjecture to a precise analytic question: does the Poisson integral $h$ with boundary values $\log|f(e^{i\varphi})|$ and $h(0)=0$ satisfy $\max_{r\in[0,1]}h(re^{i\theta})\leq 0$ for some $\theta$ with $h(e^{i\theta})\leq 0$? This connects the problem to the theory of harmonic measure and the geometry of Poisson integrals.

Problem~\#1120 is naturally part of the \emph{lemniscate cluster}: Problems \#114 (maximal lemniscate length), \#509 (covering radii), \#1038 (minimal area), and \#1120 (shortest path), all concerning extremal properties of the sublevel set $\{|f(z)|\leq 1\}$ for monic polynomials with constrained roots. Progress on any one of these illuminates the others through the shared framework of logarithmic potential theory.

\subsection*{Acknowledgments}
Computational infrastructure provided by Mendoza Lab. The Erd\H{o}s Atlas (\url{https://erdosatlas.org}) was used for problem tracking and morphism analysis.

\begin{thebibliography}{10}

\bibitem{EHP58}
P.~Erd\H{o}s, F.~Herzog, and G.~Piranian.
\newblock Metric properties of polynomials.
\newblock {\em J.\ Anal.\ Math.}, 6:125--148, 1958.

\bibitem{EP}
Erd\H{o}s Problems.
\newblock \url{https://www.erdosproblems.com/1120}.

\bibitem{FN25}
A.~Fryntov and F.~Nazarov.
\newblock On the local extremality of lemniscates of $z^n-1$.
\newblock Preprint, 2025.

\bibitem{KLR25}
M.~Krishnapur, E.~Lundberg, and K.~Ramachandran.
\newblock On the minimal area of polynomial lemniscates.
\newblock arXiv:2503.18270, 2025.

\bibitem{Ransford}
T.~Ransford.
\newblock {\em Potential Theory in the Complex Plane}.
\newblock Cambridge University Press, 1995.

\bibitem{Tao25}
T.~Tao.
\newblock The Erd\H{o}s--Herzog--Piranian conjecture for large polynomial degrees.
\newblock arXiv:2512.12455, 2025.

\end{thebibliography}

\end{document}
